\section{Threats to Validity}
\label{sec:threats_validity}

\subsection{Construct validity}
\label{sec:construct_validity}

M2 uses churn as a proxy for change magnitude, which does not directly encode semantic complexity or defect risk.
M3 cadence features are sensitive to timestamp semantics; commit-date/tag-date drift is documented and treated as a robustness concern.
Automated accounts can also inflate commit-frequency and ownership/cadence interpretations, so M3 and M4 use human-only
outputs as the primary population and keep all-account outputs for sensitivity checks.
To reduce construct ambiguity, the study keeps raw fields (added, deleted, commits, activity days) alongside derived indicators.

\subsection{Internal validity}
\label{sec:internal_validity}

Potential extraction errors were mitigated by deterministic windows, explicit cleaning rules, and arithmetic checks
(\texttt{churn} $=$ \texttt{added} $+$ \texttt{deleted}). Developer identity is canonicalized by email during extraction,
window totals are recomputed and validated from clean developer rows, and bot-sensitivity outputs validate that all-account
totals remain consistent before deriving human-only tables.

\subsection{External validity}
\label{sec:external_validity}

The Initial Cohort includes four mature Python projects only, so generalization to other ecosystems is limited.
This is intentional: Phase I focuses on pipeline correctness; the Extended Cohort adds six more repositories to test transferability
across a larger and more heterogeneous sample.

\subsection{Conclusion validity}
\label{sec:conclusion_validity}

No inferential claims are made in Phase I. Comparative conclusions are deferred to Phase II and will be supported by
explicit stability/ranking analyses and the predefined feasibility rubric. This sequencing reduces risk of premature
metric selection bias.






